%% Chapitre 14 — Cheat-sheet : commandes et raccourcis essentiels
%% RÈGLE : uniquement des commandes sémantiques bookforge.

\chapitre{Cheat-sheet : l'essentiel en un coup d'œil}

\introChapitre{
  Toutes les commandes, raccourcis et patterns utiles au quotidien, regroupés
  pour une consultation rapide. À garder sous la main lors de tes premières
  semaines avec Claude Code.
}

\section{Démarrage et navigation}

\paragraphe{Les commandes de base pour lancer et piloter une session :}

\begin{extraitcode}{bash}
# Lancer Claude Code dans le répertoire courant
claude

# Lancer avec une tâche directe (mode non-interactif)
claude -p "Décris la structure du projet"

# Reprendre la dernière session
claude --continue

# Reprendre une session spécifique
claude --resume
\end{extraitcode}

\begin{extraitcode}{bash}
# Dans une session : commandes slash essentielles
/help          # Liste toutes les commandes disponibles
/clear         # Vider le contexte de la session en cours
/compact       # Résumer et compresser l'historique
/memory        # Afficher ou éditer la mémoire projet
/cost          # Afficher le coût de la session en cours
/exit          # Quitter proprement
\end{extraitcode}

\section{Gérer les permissions}

\paragraphe{Contrôler ce que Claude Code peut faire sans te demander :}

\begin{extraitcode}{bash}
# Autoriser une commande spécifique pour la session
/allowed-tools

# Lancer en mode lecture seule (aucune modification)
claude --no-write

# Lancer en mode autonome (approuve tout automatiquement)
# À réserver aux pipelines CI, jamais en interactif
claude --dangerously-skip-permissions
\end{extraitcode}

\begin{avertissement}
\code{--dangerously-skip-permissions} supprime toutes les confirmations. Ne
l'utilise qu'en environnement isolé (conteneur, CI) sur du code que tu peux
jeter. Jamais sur ta machine de développement principale.
\end{avertissement}

\section{Fichier CLAUDE.md — mémoire projet}

\paragraphe{Le fichier \code{CLAUDE.md} à la racine du dépôt est chargé
automatiquement à chaque session. Structure recommandée :}

\begin{extraitcode}{text}
# Conventions du projet

## Commandes essentielles
- Tests : `make test` ou `npm test`
- Lint : `make lint`
- Build : `make build`

## Règles absolues
- Ne jamais modifier `generated/` (fichiers auto-produits)
- Toujours tester avant de proposer un diff
- Nouvelle dépendance = justification obligatoire

## Architecture
- Logique métier dans `src/core/`
- Handlers HTTP dans `src/api/`
- Tests dans `tests/` (miroir de `src/`)
\end{extraitcode}

\begin{astuce}
Garde le \code{CLAUDE.md} sous 200 lignes. Au-delà, il occupe trop de fenêtre
de contexte pour chaque session. L'essentiel seulement : ce qui change le
comportement de l'agent par rapport à ses défauts.
\end{astuce}

\section{Patterns de délégation efficaces}

\paragraphe{Formulations qui donnent de bons résultats systématiquement :}

\begin{extraitcode}{text}
# Pattern 1 — Bug précis
"Dans `src/auth/login.ts` ligne 47, la fonction `validateToken`
retourne toujours true. Message d'erreur : [coller]. Corrige uniquement
cette fonction, ne touche pas aux tests."

# Pattern 2 — Nouvelle fonctionnalité bornée
"Ajoute une fonction `formatDate(date: Date): string` dans
`src/utils/date.ts`. Format attendu : 'DD/MM/YYYY'. Ajoute un test
unitaire dans `tests/utils/date.test.ts`."

# Pattern 3 — Refactorisation contrôlée
"Extrais la logique de validation de `src/api/users.ts` dans un
fichier `src/core/validation.ts`. Ne change pas le comportement
observable, uniquement la structure."
\end{extraitcode}

\section{Vérification du travail — réflexes}

\paragraphe{À faire systématiquement après chaque délégation :}

\begin{liste}
  \element{Lire le diff proposé avant d'accepter — \code{git diff} ou la
  prévisualisation intégrée.}
  \element{Lancer les tests : une modification qui casse les tests est à
  refuser, même si « elle a l'air bonne ».}
  \element{Vérifier que le périmètre n'a pas débordé : l'agent n'a modifié
  que ce qu'on lui avait demandé.}
  \element{Pour les modifications de logique, tester le cas nominal \emphase{et}
  les cas limites à la main.}
\end{liste}

\section{Hooks utiles}

\paragraphe{Exemples de hooks à placer dans \code{.claude/settings.json} :}

\begin{extraitcode}{json}
{
  "hooks": {
    "PreToolUse": [{
      "matcher": "Bash",
      "hooks": [{
        "type": "command",
        "command": "echo '[Hook] Commande : $CLAUDE_TOOL_INPUT_COMMAND'"
      }]
    }],
    "PostToolUse": [{
      "matcher": "Edit",
      "hooks": [{
        "type": "command",
        "command": "npm run lint --silent 2>&1 | head -5"
      }]
    }]
  }
}
\end{extraitcode}

\section{Commandes MCP courantes}

\begin{extraitcode}{bash}
# Lister les serveurs MCP connectés
claude mcp list

# Ajouter un serveur MCP local
claude mcp add mon-outil -- node /chemin/vers/server.js

# Ajouter un serveur MCP distant (SSE)
claude mcp add mon-api --transport sse https://api.exemple.com/mcp
\end{extraitcode}

\section{Surveillance des coûts}

\begin{extraitcode}{bash}
# Coût de la session courante (dans la session)
/cost

# Voir l'usage via l'API Anthropic
# → console.anthropic.com/usage
\end{extraitcode}

\begin{note}
Les sessions longues avec beaucoup de fichiers lus coûtent plus cher que les
sessions courtes et ciblées. Si une session dépasse 15 minutes sur une seule
tâche, c'est souvent le signe qu'il faut la découper.
\end{note}

\resumeChapitre{
  \element{\important{Démarrage} : \code{claude} pour une session interactive,
  \code{claude -p "..."} pour une tâche directe, \code{claude --continue} pour
  reprendre.}
  \element{\important{Contexte} : \code{/clear} vide la session, \code{/compact}
  compresse, \code{CLAUDE.md} persiste les conventions du projet.}
  \element{\important{Permissions} : configurer dans \code{.claude/settings.json}
  ce qui s'exécute sans confirmation ; \code{--dangerously-skip-permissions} pour
  la CI uniquement.}
  \element{\important{Vérification} : lire le diff, lancer les tests, contrôler
  le périmètre — à chaque délégation, sans exception.}
  \element{\important{Coûts} : \code{/cost} en session, sessions courtes et
  ciblées pour maîtriser la facture.}
}
